\section{Módulo de Inventario}

El módulo de Inventario es el núcleo operativo de la empresa, diseñado para eliminar los errores del conteo manual y garantizar la precisión de las existencias (camisetas, buzos, balones, trofeos) en tiempo real.

\subsection{Catálogo de Productos}

Desde el menú de \textit{Productos}, podrá visualizar todos los artículos comercializados por la empresa, junto con su precio de venta, costo y la cantidad a mano disponible en la tienda física de Arequipa.

\vspace{1cm}
\begin{center}
\textit{[Espacio reservado para la captura del Catálogo de Productos]}
\end{center}
\vspace{1cm}

\subsection{Actualización y Control de Stock}

Para mantener la exactitud del inventario:
\begin{itemize}
    \item \textbf{Ingresos de mercadería:} Las nuevas compras a proveedores deben registrarse a través de recepciones para incrementar el stock.
    \item \textbf{Salidas automáticas:} Cada vez que se procese un pago en el Punto de Venta, el sistema descontará automáticamente los productos vendidos.
    \item \textbf{Ajustes de inventario:} Excepcionalmente, si se detectan diferencias durante las auditorías físicas semestrales, el usuario autorizado podrá realizar un ajuste manual para sincronizar el stock físico con el digital.
\end{itemize}
